% =============================================================================
% SECTION 4 — Benchmark Factor Models
% =============================================================================


% ── SEMAPHORE ──────────────────────────────────────────────────────────────────

This section estimates three canonical benchmark spanning
regressions from the fixed-income arbitrage
literature---\citet{duarte2007risk}, \citet{fung2004hedge}, and
\citet{brooks2020active}---and tests whether the alpha of the three
strategies survives once their systematic factor exposures are
controlled for. A central question, motivated by the slow-moving
capital framework of \citet{duffie2010presidential}, is whether the
residual alpha varies systematically with intermediary funding
stress. We describe in Section~\ref{sec:bench_specs} the three
specifications and the construction of their euro-denominated
counterparts, and present in Section~\ref{sec:bench_results} the
regression results, the cross-model alpha synthesis, the regime
decomposition, the conditional alpha model of
\citet{christopherson1998conditioning}, and an out-of-sample
re-estimation that forms the benchmark hedge from past-only loadings.


% ── 4.1 BENCHMARK SPECIFICATIONS ──────────────────────────────────────────────

\subsection{Benchmark Specifications}
\label{sec:bench_specs}

Three benchmark specifications are widely used in the empirical
literature on fixed-income arbitrage and active fixed-income management,
and they provide the natural reference points against which our
strategies should be evaluated.

\citet{duarte2007risk}, in the paper that motivates our research
design, test whether fixed-income arbitrage excess returns are
compensation for exposures to equity, banking-sector, and
term-structure risk:
\begin{equation}\label{eq:duarte}
\begin{split}
  r_{i,t} \;=\; \alpha_i
  &+ \beta_1\,\textit{Mkt}_t + \beta_2\,\textit{SMB}_t
   + \beta_3\,\textit{HML}_t + \beta_4\,\textit{UMD}_t + \beta_5\,\textit{RS}_t + \beta_6\,\textit{RI}_t\\
  &+ \beta_7\,\textit{RB}_t + \beta_8\,\textit{R2}_t + \beta_9\,\textit{R5}_t
   + \beta_{10}\,\textit{R10}_t + \varepsilon_{i,t},
\end{split}
\end{equation}
where $\textit{Mkt}$, $\textit{SMB}$, $\textit{HML}$, and
$\textit{UMD}$ are the Fama--French--Carhart equity factors (market, size, value, and momentum),
$\textit{RS}$ is the S\&P bank stock index excess return,
$\textit{RI}$ and $\textit{RB}$ are A/BAA-rated industrial and bank
bond index excess returns, and $\textit{R2}$, $\textit{R5}$,
$\textit{R10}$ are the CRSP Fama bond portfolios at two-, five-, and
ten-year maturities.

\citet{fung2004hedge} propose a seven-factor \emph{asset-based
style} model that tests whether hedge fund returns are explained by
equity, fixed-income, and trend-following exposures:
\begin{equation}\label{eq:funghsieh}
\begin{split}
  r_{i,t} \;=\; \alpha_i
  &+ \beta_1\,\textit{S\&P}_t + \beta_2\,\textit{SC\text{-}LC}_t
   + \beta_3\,\textit{BdOpt}_t + \beta_4\,\textit{FXOpt}_t \\
  &+ \beta_5\,\textit{ComOpt}_t
   + \beta_6\,\textit{10Y}_t
   + \beta_7\,\textit{CredSpr}_t + \varepsilon_{i,t},
\end{split}
\end{equation}
where $\textit{S\&P}$ is the S\&P~500 excess return,
$\textit{SC\text{-}LC}$ is the small-minus-large-cap total return,
$\textit{BdOpt}$, $\textit{FXOpt}$, and $\textit{ComOpt}$ are the three
trend-following factors constructed as
lookback straddle returns on bond, currency, and commodity futures,
$\textit{10Y}$ is the excess return on a 10-year government bond
total-return index, and $\textit{CredSpr}$ is the return spread between
BAA-rated corporate bonds and 10-year government bonds. The last two
factors replace the yield and spread \emph{changes} of the original
specification with their tradable total-return counterparts: a yield
change is not the return on an investable position, so the intercept of
a regression on yield changes lacks the Jensen-alpha interpretation;
with all seven regressors entered as realised excess returns, the
intercept retains its meaning as abnormal return.

\citet{brooks2020active} take a different angle,
decomposing the returns of active fixed-income managers into eight
\emph{traditional risk premia} organised along three channels of
risk-taking---duration, credit, and volatility selling:
\begin{equation}\label{eq:brooks}
\begin{split}
  r_{i,t} \;=\; \alpha_i
  &+ \beta_1\,\textit{Term}_t + \beta_2\,\textit{GlTerm}_t
   + \beta_3\,\textit{GlAgg}_t + \beta_4\,\textit{InflLink}_t 
   + \beta_5\,\textit{CorpCred}_t \\
  &+ \beta_6\,\textit{EmDebt}_t + \beta_7\,\textit{EmCurr}_t 
  + \beta_8\,\textit{USTVol}_t + \varepsilon_{i,t},
\end{split}
\end{equation}
where $\textit{Term}$ and $\textit{GlTerm}$ are the total returns on
the US Treasury and the currency-hedged global Treasury indices, both
in excess of cash (3-month Treasury bills), $\textit{GlAgg}$ is the
total return on the currency-hedged global aggregate bond index in
excess of cash, $\textit{InflLink}$ is the total return on the
currency-hedged global inflation-linked Treasury index in excess of
cash, $\textit{CorpCred}$ is a 50/50
blend of a US high-yield corporate bond index (in excess of
duration-matched Treasuries) and a leveraged loan index (in
excess of 3-month LIBOR), $\textit{EmDebt}$ is the duration-adjusted
excess return on an emerging-market debt index,
$\textit{EmCurr}$ is the return on an equal-weighted basket of
emerging-market currencies, and $\textit{USTVol}$ is the return on
delta-hedged straddles on 10-year Treasury futures.

Each benchmark is estimated in two parallel panels. The first uses the
original U.S.-denominated factors as specified by the respective authors,
which is the natural way to replicate each framework as it was first
published. The second uses euro equivalents, since all three strategies
trade European instruments and may therefore be better explained by
European factor exposures. The euro panel is the primary specification
and is reported in the main text; the US-factor estimates are reported
in Appendix~\ref{app:benchmark_us}.

For the euro panel, the Fama--French factors---used directly in
\citet{duarte2007risk} and indirectly in
\citet{fung2004hedge}---are sourced from the Kenneth French data
library; because these are reported in USD, we convert them to EUR
following \citet{gluck2021currency}. For long factors (market excess
return):
\begin{equation}\label{eq:gluck_long}
  F_{t}^{\textsc{eur}}
  \;=\; \frac{1 + F_{t}^{\textsc{usd}} + r_{f,t}^{\textsc{usd}}}
             {1 + r_{t}^{\textsc{usd/eur}}}
        \;-\; 1 \;-\; r_{f,t}^{\textsc{eur}},
\end{equation}
and for long-short factors (SMB, HML, UMD):
\begin{equation}\label{eq:gluck_ls}
  \mathit{LS}_{t}^{\textsc{eur}}
  \;=\; \frac{\mathit{LS}_{t}^{\textsc{usd}}}
             {1 + r_{t}^{\textsc{usd/eur}}},
\end{equation}
where $F_{t}^{\textsc{usd}}$ and $\mathit{LS}_{t}^{\textsc{usd}}$ are
the original factors in US dollars, $r_{f,t}^{\textsc{usd}}$ and
$r_{f,t}^{\textsc{eur}}$ are the US and euro risk-free rates
(Fama--French T-bill rate and one-month Euribor, respectively), and
$r_{t}^{\textsc{usd/eur}}$ is the period return on the dollar-to-euro
exchange rate.
All remaining factors enter the euro panel with native European
counterparts: European bank stocks, euro-area corporate bond indices,
and German government bonds in the Duarte specification; euro-area
term, credit, and inflation-linked premia in Brooks et al.; and the
10-year Bund yield and Baa--Bund credit spread in Fung and Hsieh.
The Fung and Hsieh trend-following lookback straddle factors are
currency-independent by construction (D.~Hsieh, private
communication) and enter both panels identically. The complete factor
list is reported in Appendix~\ref{app:factor_list}.

All specifications are estimated by ordinary least squares with
Newey--West HAC standard errors \citep{newey1987simple}, with the
truncation lag set by the automatic rule of \citet{newey1994automatic},
$\lfloor 4\,(T/100)^{2/9} \rfloor$, which equals four lags for every
sample in this paper. Variance inflation factors for each specification are reported in
Appendix~\ref{app:vif}. All factors exhibit VIF below~$10$ except
for the Fama bond portfolios (R2, R5, R10) in the Duarte
specification (VIF~$\approx 15$--$55$), which reflects the
well-known correlation structure of maturity-sorted Treasury
returns.


% ── 4.2 RESULTS AND DISCUSSION ────────────────────────────────────────────────

\subsection{Results and Discussion}
\label{sec:bench_results}

\paragraph{Factor Model Regressions.}

Tables~\ref{tab:duarte}, \ref{tab:fung_hsieh}, and~\ref{tab:active_fi}
report the full euro-panel regression results; the US-factor
counterparts are Tables~\ref{tab:duarte_us}, \ref{tab:fung_hsieh_us},
and~\ref{tab:active_fi_us} in Appendix~\ref{app:benchmark_us}. Across
all three frameworks, all three strategies, and both currency
denominations, the estimated intercept is positive and statistically
significant at the one-percent level. The annualised alphas are economically meaningful:
approximately $+5\%$ p.a.\ for the CDS--Bond basis across all
specifications, $+1.6$--$3.5\%$ for iTraxx Skew, and
$+1.6$--$2\%$ for BTP~Italia.

The adjusted $R^2$ is low across all specifications, ranging from one
to eighteen percent, indicating that the standard factor sets leave
most of the return variation unexplained---a question we take up in
Section~\ref{sec:factor_selection}. The significant factor loadings that do emerge differ across
strategies and are economically intuitive. The CDS--Bond basis loads
on credit and rates factors: the A-rated industrial corporate bond
return is positive and significant in the Duarte specification in both
currency denominations, and the change in the 10-year Treasury yield
enters negatively in the Fung--Hsieh specification, again in both
denominations---the natural exposures of a long-credit, cash-financed
position. BTP~Italia and iTraxx Skew load instead on equity
factors: the market, size, and value factors in Duarte and the
small-minus-large-cap spread in Fung--Hsieh, with the broad equity
market significant for BTP~Italia but not for iTraxx Skew. These
loadings carry negative signs, consistent with the modest indirect
equity exposure embedded in the two basis trades. The Brooks et al.\ specification,
which uses broader risk premia rather than individual asset returns,
produces fewer significant loadings but delivers comparable alpha estimates,
reinforcing that the intercepts are not an artefact of any single factor set.

\paragraph{Cross-Model Comparison and Regime Analysis.}

Table~\ref{tab:alpha_synthesis} consolidates the alpha estimates from
the three frameworks. Panel~A
reports the unconditional alphas side by side: for each strategy, all
three frameworks deliver positive and significant alpha, and the
estimates are consistent across specifications despite their different
factor structures.

The discrete regime decomposition---alpha estimated separately in the
LOW, MEDIUM, and HIGH stress regimes defined on the level of the iTraxx
Europe Main 5-year spread, with the same classification introduced in
Section~\ref{sec:performance} (LOW below $60$~basis points, MEDIUM
between $60$ and $100$, HIGH above $100$)---is reported for each
framework in Appendix~\ref{app:subperiod}. In every cell of those
tables, alpha rises monotonically from the LOW to the HIGH regime, and
the HIGH-regime alpha is a multiple of the LOW-regime alpha---in most
specifications by a factor of three to four. The pattern holds across
all three benchmark specifications and across all three strategies. The
regime cutoffs are not arbitrary: they are the same 60 and 100
basis-point thresholds that govern the stress-dependent transaction-cost
tiers applied to every return series in the paper
(Appendix~\ref{app:fees}); and the conditional model in Panel~B, which
uses the full continuous variation in the stress proxy rather than a
discrete split, confirms that the regime effect does not hinge on the
choice of cutoffs.

This regime dependence is consistent with the theoretical mechanism that
\citet{duffie2010presidential} formalises in his presidential address.
When intermediary capital is abundant, arbitrage forces are strong and
mispricings are quickly closed, compressing the residual premium that
remains after compensating for the systematic factor loadings. When
intermediary capital is scarce---as it is during episodes of widening
credit spreads---the arbitrageurs who ordinarily close the gap retreat
from the trade or are forced to unwind existing positions, the
mispricings widen, and the residual premium for those who can still
deploy capital rises sharply. This is precisely the regime dependence
documented in those tables, and it is also the empirical pattern that
motivates the cross-strategy analysis in
Section~\ref{sec:interdependencies}, where we test whether the same
intermediary balance-sheet shocks drive co-movement of mispricings across
the three strategies.

\paragraph{Conditional Alpha.}

Panel~B of Table~\ref{tab:alpha_synthesis} formalises the regime effect
using the conditional performance evaluation framework of
\citet{christopherson1998conditioning}, which lets the alpha intercept vary
continuously with a publicly observable state variable rather than across
a binary regime split. The specification is
\begin{equation}\label{eq:ferson_schadt}
  r_{i,t} \;=\; \alpha_0 \;+\; \alpha_1\, z_t
  \;+\; \boldsymbol{\beta}'\, \mathbf{X}_{t} \;+\; \varepsilon_{t},
\end{equation}
where $r_{i,t}$ is the strategy excess return, $\mathbf{X}_t$ collects the factor
returns of the chosen benchmark, and $z_t$ is the standardised level of
the iTraxx Europe Main 5-year spread:
$z_t = (s_t - \bar{s})/\hat{\sigma}_s$. The state variable $z_t$ is
constructed to be mean-zero and unit-variance over the sample, so the
intercept $\alpha_0$ retains its interpretation as the unconditional alpha
and the slope $\alpha_1$ measures the additional alpha earned per
one-standard-deviation increase in funding stress. Standard errors are
again computed with the \citet{newey1994automatic} HAC correction. 
The \citet{christopherson1998conditioning} formulation improves on the binary
split in Panel~B by using the full continuous variation in the stress
proxy rather than collapsing it into a dummy, which is more efficient
when the underlying state varies meaningfully over time.

The estimates of $\alpha_1$ in Panel~B are positive for all three
strategies and all three benchmark specifications, and the
unconditional intercept $\alpha_0$ remains significant in every case,
confirming that a baseline alpha persists even after conditioning on
the stress state. For CDS--Bond Basis and iTraxx Skew, $\alpha_1$
is significant at the one-percent level across all three frameworks: a
one-standard-deviation increase in the iTraxx Main spread raises
CDS--Bond alpha by approximately three percentage points per year and
iTraxx alpha by one to three percentage points. For BTP~Italia, $\alpha_1$ is positive in all specifications but
less precisely estimated (significant only at the 5--10\% level),
consistent with the shorter sample available for this strategy. For all three strategies, the implied alpha at
$z = +1\sigma$ is approximately three to nine times larger than at
$z = -1\sigma$, indicating that the bulk of the unconditional alpha
is earned during periods of elevated funding stress. Reading the regime
decomposition and Panel~B together, the unconditional alpha masks a pronounced time variation
consistent with the slow-moving capital mechanism of
\citet{duffie2010presidential}.

\paragraph{Out-of-Sample Alpha.}

A residual concern, central to the literature on out-of-sample
performance evaluation \citep{moreira2017volatility, cederburg2020performance},
is that the benchmark loadings in
Tables~\ref{tab:duarte}--\ref{tab:active_fi} are themselves estimated in
sample, so the surviving alpha could in principle reflect an
in-sample-optimal combination of the factors rather than an abnormal
return that an investor could have earned in real time. Panel~C of
Table~\ref{tab:alpha_synthesis} addresses this directly. For each
strategy and each benchmark, we re-estimate the factor loadings on an
expanding past-only window in the recursive out-of-sample design of
\citet{welch2008comprehensive}, form the factor hedge from the lagged
loadings, and define the out-of-sample alpha as the mean of the resulting hedged return,
tested with the \citet{newey1994automatic} HAC correction. The expanding
window is initialised with a 60-month minimum, so that the first
benchmark fit has at least five observations per parameter under the
ten-factor \citet{duarte2007risk} specification. The out-of-sample alpha remains
positive and statistically significant for all three strategies and all
three frameworks. The point estimates are broadly comparable to, and in
most cases modestly below, their in-sample counterparts in
Panel~A---the attenuation expected once the hedge can no longer be
fitted to the same returns it is applied to---and the ranking across
strategies is preserved, with the CDS--Bond basis the largest and
BTP~Italia, estimated over the shortest sample, the least precisely
identified. The rolling-window counterpart, which additionally relaxes
the assumption of constant loadings, is reported in
Appendix~\ref{app:oos_rolling} and delivers the same conclusion.

\paragraph{Rolling Alpha.}

A natural concern with regime-based or conditional analyses is that the
estimated relationship between alpha and stress may be driven by a
single sub-period---perhaps the global financial crisis of 2008--2009,
or the European sovereign debt crisis of 2011--2012, or the COVID-19
shock of March 2020---rather than by a stable structural feature of the
data. Figure~\ref{fig:rolling_alpha_composite} addresses this concern
directly by plotting rolling 36-month alpha estimates for each strategy
under each of the three benchmark specifications. The rolling alpha is
positive across essentially the whole sample for all three strategies,
so the estimates are not the product of a single favourable sub-period
but persist throughout. The three benchmark estimates move broadly
together within each panel, reinforcing the cross-model robustness
documented in Panel~A of Table~\ref{tab:alpha_synthesis}. Detailed
sub-period splits and regime decompositions for each framework
individually are reported in Appendix~\ref{app:subperiod}.

Taken together, the five sets of results in this section establish two
findings that are central to the rest of the paper. First, the alpha of
the three arbitrage strategies is robust: it survives controls based on
three widely used benchmark frameworks for fixed-income arbitrage, in
both currency denominations; it is not an artefact of in-sample factor
combination, surviving when each benchmark hedge is re-estimated out of
sample; and it does so for every cross-section and every sub-period that
we examine. Second, the alpha is not constant. It
rises systematically when intermediary capital is scarce, consistent with the
slow-moving capital theory, and the regime effect is large: in several
specifications the alpha earned in high-stress months is a multiple of
the alpha earned in normal months. The first finding indicates that the observed excess returns are not
explained by the systematic exposures that the prior literature has
identified. The second finding points to a deeper economic mechanism,
but it does not yet identify which specific risk factors are
responsible for the residual variation. We turn to that question in
Section~\ref{sec:factor_selection}, where we use penalised regression
methods on a much larger candidate factor universe to identify the
factors that these three benchmarks omit.